\documentclass[a3paper, landscape]{article}
\usepackage[left=1cm, right=1cm, top=1cm, bottom=1cm]{geometry}
\usepackage[T2A]{fontenc}
\usepackage[utf8]{inputenc}
\usepackage[russian]{babel}
\thispagestyle{empty}

\usepackage{tikz}
\usetikzlibrary{positioning, calc, matrix}

\tikzset{
    tbl/.style={
        matrix of nodes,
        draw, thick,
        inner sep=0pt,
        column sep=0pt,
        row sep=0pt,
        nodes={
            anchor=center,
            font=\ttfamily\scriptsize,
            inner xsep=4pt,
            inner ysep=3pt,
            minimum width=14em,
            minimum height=1.6em,
            draw, thin,
            align=left,
        },
        row 1/.style={
            nodes={
                font=\bfseries\small,
                minimum height=2em,
                align=center,
                fill=white,
            },
        },
    },
    ann/.style={font=\scriptsize},
    lnk/.style={thick},
    crd/.style={font=\scriptsize, fill=white, inner sep=1pt},
}

\begin{document}
\begin{figure}[p]
\centering
\begin{tikzpicture}[node distance=3em]

% === Ряд 1: Издательство — Гарнитура — Категория ===

\matrix[tbl] (pub) at (0,0) {
    Издательство \\
    id\_Издательство \\
    Название \\
    Страна \\
    Сайт \\
    Год\_основания \\
    Email \\
};
\node[ann, above=0.3em of pub] {--; --; --};

\matrix[tbl] (ff) at (25em,0) {
    Гарнитура \\
    id\_Гарнитура \\
    Название \\
    Год\_создания \\
    Описание \\
    Число\_начертаний \\
    id\_Издательство \\
    id\_Категория \\
};
\node[ann, above=0.3em of ff] {--; --; --};

\matrix[tbl] (cat) at (50em,0) {
    Категория \\
    id\_Категория \\
    Название \\
    Описание \\
    Иконка \\
    Число\_гарнитур \\
};
\node[ann, above=0.3em of cat] {--; --; --};

% === Ряд 2: Лицензия — Начертание — Формат ===

\matrix[tbl] (lic) at (0,-18em) {
    Лицензия \\
    id\_Лицензия \\
    Тип \\
    Срок\_действия \\
    Макс\_число\_станций \\
    Стоимость \\
    Дата\_приобретения \\
    id\_Издательство \\
    id\_Гарнитура \\
};
\node[ann, above=0.3em of lic] {--; --; --};

\matrix[tbl] (tf) at (25em,-18em) {
    Начертание \\
    id\_Начертание \\
    Название \\
    Насыщенность \\
    Наклон \\
    Размер\_файла \\
    id\_Гарнитура \\
    id\_Формат \\
};
\node[ann, above=0.3em of tf] {--; --; --};

\matrix[tbl] (fmt) at (50em,-18em) {
    Формат \\
    id\_Формат \\
    Название \\
    Расширение \\
    Год\_появления \\
    Описание \\

    Поддержка\_веба \\
};
\node[ann, above=0.3em of fmt] {--; --; --};

% === Ряд 3: Язык — Алфавит — Установка — Раб. станция — ОС ===

\matrix[tbl] (lang) at (0,-54em) {
    Язык \\
    id\_Язык \\
    Название \\
    Код\_ISO \\
    Направление\_письма \\
    Тип\_письменности \\
};
\node[ann, above=0.3em of lang] {--; --; --};

\matrix[tbl] (alph) at (16em,-36em) {
    Алфавит \\
    id\_Алфавит \\
    Число\_начерт\_в\_сист. \\
    Число\_поддерж\_гарн. \\
    Дата\_проверки \\
    Статус \\
    id\_Начертание \\
    id\_Язык \\
};
\node[ann, above=0.3em of alph] {--; --; --};

\matrix[tbl] (inst) at (40em,-36em) {
    Установка \\
    id\_Установка \\
    Дата\_установки \\
    Дата\_удаления \\
    Путь\_к\_файлу \\
    Статус \\
    Дата\_последнего\_исп. \\
    id\_Начертание \\
    id\_Рабочая\_станция \\
};
\node[ann, above=0.3em of inst] {--; --; --};

\matrix[tbl] (ws) at (64em,-36em) {
    Рабочая станция \\
    id\_Рабочая\_станция \\
    Инвентарный\_номер \\
    Процессор \\
    Объём\_памяти \\
    Тип\_накопителя \\
    Графический\_адаптер \\
    id\_ОС \\
};
\node[ann, above=0.3em of ws] {--; --; --};

\matrix[tbl] (os) at (64em,-54em) {
    ОС \\
    id\_ОС \\
    Название \\
    Версия \\
    Разрядность \\
    Год\_выпуска \\
};
\node[ann, above=0.3em of os] {--; --; --};

% ==================== СВЯЗИ ====================

% --- Ряд 1: горизонтальные ---

% Гарнитура.id_Издательство → Издательство.id_Издательство
\draw[lnk] (ff-7-1.west) -- node[crd, above] {$\infty$} ++(-1.5em,0)
    -- node[crd, above, pos=0.5] {} ([xshift=1.5em]pub-2-1.east)
    -- node[crd, above] {1} (pub-2-1.east);

% Гарнитура.id_Категория → Категория.id_Категория
\draw[lnk] (ff-8-1.east) -- node[crd, above] {$\infty$} ++(1.5em,0)
    -- ([xshift=-1.5em]cat-2-1.west)
    -- node[crd, above] {1} (cat-2-1.west);

% --- Ряд 1→2: вертикальные ---

% Гарнитура.id_Гарнитура → Начертание.id_Гарнитура
\draw[lnk] (ff-2-1.east) -- ++(2em,0) node[crd, above] {1}
    |- node[crd, pos=0.75, above] {$\infty$} (tf-7-1.east);

% Издательство.id_Издательство → Лицензия.id_Издательство
\draw[lnk] (pub-2-1.west) -- ++(-1.5em,0) node[crd, above] {1}
    |- node[crd, pos=0.75, above] {$\infty$} (lic-8-1.west);

% Гарнитура.id_Гарнитура → Лицензия.id_Гарнитура
\draw[lnk] (ff-2-1.west) -- ++(-3em,0) node[crd, above] {1}
    |- node[crd, pos=0.75, above] {$\infty$} (lic-9-1.east);

% --- Ряд 2: горизонтальные ---

% Начертание.id_Формат → Формат.id_Формат
\draw[lnk] (tf-8-1.east) -- node[crd, above] {$\infty$} ++(1.5em,0)
    -- ([xshift=-1.5em]fmt-2-1.west)
    -- node[crd, above] {1} (fmt-2-1.west);

% --- Ряд 2→3: вертикальные ---

% Начертание.id_Начертание → Алфавит.id_Начертание
\draw[lnk] (tf-2-1.west) -- ++(-2em,0) node[crd, above] {1}
    |- node[crd, pos=0.75, above] {$\infty$} (alph-7-1.west);

% Начертание.id_Начертание → Установка.id_Начертание
\draw[lnk] (tf-2-1.west) -- ++(-2em,0) node[crd, above] {1}
    |- ++(0,-14em)
    |- ([yshift=2em]inst-8-1.west)
    -- node[crd, left] {$\infty$} (inst-8-1.west);

% --- Ряд 3: горизонтальные ---

% Алфавит.id_Язык → Язык.id_Язык
\draw[lnk] (alph-8-1.west) -- node[crd, above] {$\infty$} ++(-1.5em,0)
    -- ([xshift=1.5em]lang-2-1.east)
    -- node[crd, above] {1} (lang-2-1.east);

% Установка.id_Рабочая_станция → Рабочая_станция.id_Рабочая_станция
\draw[lnk] (inst-9-1.east) -- node[crd, above] {$\infty$} ++(8em,0)
    |- ([xshift=-1.5em]ws-2-1.west)
    |- node[crd, above] {1} (ws-2-1.west);

% Рабочая_станция.id_ОС → ОС.id_ОС
\draw[lnk] (ws-8-1.east) -- ++(2em,0) node[crd, above] {$\infty$}
    |- node[crd, left, pos=0.75] {1} (os-2-1.east);

\end{tikzpicture}
\caption{Схема базы данных на русском языке}
\end{figure}
\end{document}
